Fourteen states currently use some form of independent or bipartisan commission to draw congressional or state legislative districts. None currently requires algorithmic redistricting. This paper analyzes three modes for integrating the \texttt{bisect} recursive bisection algorithm into commission redistricting processes: Mode~(a), algorithm-as-sole-mapmaker (the commission approves or rejects the algorithmic output); Mode~(b), algorithm-as-first-draft (the algorithm produces a baseline that the commission modifies under documented criteria); and Mode~(c), algorithm-as-audit (the commission draws maps independently, then uses the algorithm to verify compliance with stated criteria). We find that Mode~(b) is most feasible for existing commissions: it preserves the commission's deliberative function, provides a defensible neutral starting point, and creates a verifiable record of departures. We develop the Iowa model as the strongest existing analogue to algorithmic commission redistricting: Iowa's Legislative Services Agency (LSA) draws maps using strictly enumerated criteria with no partisan data, and the legislature votes up or down on the LSA's proposals without amendments. The \texttt{bisect} platform extends the Iowa model by replacing the LSA's manual map-drawing with deterministic algorithmic output, eliminating the residual human manipulation space in criteria-only approaches. We provide a detailed Commission Adjustment Protocol specifying the procedures, documentation requirements, and certification steps for Mode~(b) integration.
